\documentclass[preprint,12pt,3p]{elsarticle}

\usepackage{amssymb}
\usepackage{amsmath}
\usepackage{comment}
\usepackage{natbib}
\usepackage{xcolor}
\usepackage{enumitem}
\usepackage{multirow}
\bibliographystyle{plainnat}
\usepackage{booktabs}
\usepackage{subcaption}
\usepackage[normalem]{ulem}
\usepackage{array}
\usepackage{hyperref}
\hypersetup{colorlinks=true,linkcolor=black,citecolor=black,urlcolor=blue}

% ── review layer ─────────────────────────────────────────────────────────────
\definecolor{annred}{rgb}{0.78,0.00,0.00}
\definecolor{annblue}{rgb}{0.00,0.22,0.85}
\definecolor{annpurple}{rgb}{0.48,0.00,0.60}
\definecolor{anngreen}{rgb}{0.00,0.40,0.13}
\definecolor{annorange}{rgb}{0.80,0.40,0.00}
\definecolor{anngray}{rgb}{0.45,0.45,0.45}
\definecolor{provyellow}{rgb}{1.00,0.94,0.35}
\definecolor{darkgreen}{rgb}{0.0,0.3,0.0}

\newcommand{\red}[1]{\textcolor{annred}{#1}}
\newcommand{\black}[1]{\textcolor{black}{#1}}
\newcommand{\orange}[1]{\textcolor{annorange}{#1}}
\newcommand{\purple}[1]{\textcolor{annpurple}{#1}}
\newcommand{\green}[1]{{\color{darkgreen}#1}}

\newcommand{\CMT}[1]{{\color{anngreen}\textsf{\small [#1]}}}
\newcommand{\FIX}[1]{{\color{annred}#1}}
\newcommand{\TODO}[1]{{\color{annblue}\textsf{\small [TO DO: #1]}}}
\newcommand{\NEW}[1]{{\color{annpurple}#1}}
\newcommand{\CODE}[1]{{\color{annorange}\textsf{\small [CODE: #1]}}}
\newcommand{\DEL}[1]{{\color{anngray}\sout{#1}}}

% yellow = provisional, will move when the remaining runs land
\newcommand{\prov}[1]{{\setlength{\fboxsep}{1pt}\colorbox{provyellow}{#1}}}
\newcommand{\provlow}[1]{{\setlength{\fboxsep}{1pt}\colorbox{provyellow}{\color{annred}\textbf{#1}}}}

\newcommand{\prio}[1]{{\color{annred}\textbf{\textsf{#1}}}}

\newenvironment{ann}[2]{%
  \par\addvspace{4pt}%
  \begin{list}{}{\leftmargin=1.6em\rightmargin=0.3em\topsep=0pt\parsep=3pt\itemsep=0pt\listparindent=0pt}%
  \item[]\color{#1}\small\textbf{\textsf{#2}}\par\nopagebreak\vspace{1pt}\color{#1}\small\ignorespaces}
  {\end{list}\addvspace{4pt}}
\newenvironment{acmt}{\begin{ann}{anngreen}{$\blacktriangleright$\ COMMENT}}{\end{ann}}
\newenvironment{afix}{\begin{ann}{annred}{$\blacktriangleright$\ CORRECTION}}{\end{ann}}
\newenvironment{atodo}{\begin{ann}{annblue}{$\blacktriangleright$\ TO DO}}{\end{ann}}
\newenvironment{anew}{\begin{ann}{annpurple}{$\blacktriangleright$\ SUGGESTED TEXT}}{\end{ann}}
\newenvironment{acode}{\begin{ann}{annorange}{$\blacktriangleright$\ CODE vs.\ PAPER}}{\end{ann}}

\begin{comment}
\journal{Energy}
\end{comment}
\begin{document}

\begin{frontmatter}
\fntext[label2]{Corresponding author: celine.pagnier@ntnu.no}
\title{Joint scheduling of driver activities and charging for long distance electric truck operation under travel time uncertainty}

\author[inst1, label1]{C\'eline Pagnier}
\author[inst2]{Margaux Nattaf}
\author[inst2]{Marie-Laure Espinouse}
\author[inst3]{Cyrille Briand}
\author[inst1]{Steffen J.S. Bakker}

\affiliation[inst1]{organization={Department of Industrial Economics and Technology Management, Norwegian University of Science and Technology},
            addressline={Alfred Getz veg 3}, city={Trondheim}, postcode={NO-7491}, country={Norway}}
\affiliation[inst2]{organization={Univ. Grenoble Alpes, CNRS, Institute of Engineering Univ. Grenoble Alpes, G-SCOP},
            city={Grenoble}, postcode={38000}, country={France}}
\affiliation[inst3]{organization={LAAS-CNRS, Universit\'e de Toulouse, CNRS, UT3},
            city={Toulouse}, postcode={31000}, country={France}}

\begin{abstract}
Battery electric trucks are entering long-haul freight operations, but various additional constraints compared to conventional diesel trucks discourage freight companies from transitioning their fleets.
As charging infrastructure expands, long-haul routes become increasingly feasible, and efficient schedules seek to synchronize charging stops with mandatory driver breaks. However, this synchronization remains fragile due to battery limitations and operational uncertainty, and becoming stranded due to lack of battery is unacceptable.
This paper studies the joint scheduling of driver activities, such as breaks and rests, and partial non-linear fast charging activities for a battery electric truck serving a fixed sequence of customers on a long-haul route under travel-time uncertainty.

We formulate a mixed-integer linear program capturing the Hours of Service regulations, including split breaks, reduced rests, extended driving allowances, and the interaction between charging and break. Around this model, we build and compare six decision architectures that differ in when decisions are committed and what information they use: a scenario-based rolling-horizon look-ahead policy, a myopic greedy policy, a robust plan, a budgeted robust plan, a two-stage stochastic plan with online duration recourse, and a perfect-hindsight oracle.

\DEL{Experiments on synthetic corridors and a real Norwegian long-haul instance show that TO BE DONE}
\DEL{- policies that adapt schedule en route perform much better than plans that commit at departure - electrification of those corridors cost XXX in route duration compared to diesel truck - decision maker and policies should focus on investment in larger chargers in CS instead of denser network of CS, but that poses problem of grid investment for powerful chargers -}

\NEW{Experiments on \prov{1\,800} synthetic corridor instances show that policies
that adapt the schedule en route perform markedly better than plans committed at
departure: the rolling-horizon and two-stage architectures stay within
\prov{12\,\%} and \prov{15\,\%} of the hindsight optimum on average, against
\prov{25\,\%} for a budgeted robust plan and \prov{35\,\%} for a box robust plan,
and the entire penalty of committing early is paid in customer time-window
compliance rather than in driving time. Electrifying these corridors costs
\prov{8--11\,\%} of route duration against a diesel truck on identical routes and
identical travel-time realizations, rising to \prov{20--23\,\%} for an operator
scheduling myopically, so roughly half of the observed electrification penalty is
a planning failure rather than a physical cost. Finally, decision makers should
direct charging-infrastructure investment towards more powerful chargers rather
than a denser network: at the spacing already mandated by Regulation (EU)
2023/1804, charger power outweighs charger density by a factor of roughly four,
although this shifts the burden onto grid connection capacity at each site.}
\end{abstract}

\begin{acmt}
\textbf{Colour code for this build.}
{\color{anngreen}green} = comment;
{\color{annred}red} = correction, with deletions struck out in {\color{anngray}\sout{grey}};
{\color{annblue}blue} = still to be done;
{\color{annpurple}purple} = suggested new text;
{\color{annorange}orange} = the manuscript and the code disagree.
\prov{Yellow} marks every number that will move when the outstanding runs land;
\provlow{yellow + red} marks a number currently computed from fewer than twenty
realizations, which should not be quoted at all until the runs complete.

\medskip
Per your status: diesel is treated as final (no highlight), the base case as
almost final (light highlight on aggregates only), 2SP as provisional
throughout, and the look-ahead sweep as complete except MIPTAIL.
All numbers below come from the tables regenerated on 2026-08-07 at 11:25--11:47.
\end{acmt}

\begin{keyword}
Electric trucks \sep TDSP \sep travel time uncertainty \sep rolling horizon
\end{keyword}
\end{frontmatter}

\section{Introduction} \label{intro}
Road freight accounts for roughly a quarter of transport $CO_2$ emissions in the European Union, and battery electric heavy-duty trucks (BEHDTs) are the leading decarbonization pathway for the segment. Their deployment in long-haul  service, however, is first constrained by battery capacity and charging infrastructure, but also by an operational factor that is absent for passenger electric vehicles: driver's regulations. Commercial drivers in the EU operate under Regulation (EC) No 561/2006 and Directive 2002/15/EC, which impose mandatory breaks after at most 4.5 h of driving, daily rest periods, and limits on shift driving time.
Because a mandatory 45-minute break is long enough for a substantial fast charge, synchronization of charging and drivers constraints is critical. Whether electric long-haul operation can match diesel route durations therefore hinges on how well charging events can be synchronized with legally required interruptions.

However, this synchronization is fragile. Travel times on long-haul corridors are uncertain, and a delay on a single leg propagates through the entire schedule: it consumes driving-time budget before the planned break location is reached, shifts arrivals outside customer service windows, and, because energy consumption depends on realized speed, perturbs the state of charge with which the truck reaches the next charging opportunity. \NEW{The two risks pull in opposite directions: a \emph{slow} leg threatens Hours-of-Service compliance and the customer window, whereas a \emph{fast} leg is the one that threatens the state of charge, since the consumption rate rises with speed over the whole operating band of a heavy vehicle.} A pre-determined schedule that nests charging inside breaks under average travel times may, in execution, force the driver into an unacceptable event such as running out of battery if the energy consumption deviates from anticipated values.

\begin{acmt}
One sentence added, because the opposite-tails structure is the justification for
the two-sided budget in Section~\ref{sec:robu} and is currently never stated.
See the orange note in Section~\ref{sec:uncertaintylit} for the numbers.
\end{acmt}

This paper studies the joint scheduling of driver activities and charging for a battery electric truck on a fixed long-haul route with multiple customer stops, under stochastic travel times. Given the sequence of customers, the decisions are where and for how long to charge (with a nonlinear, piecewise-linearized charging curve and partial charging), where to take breaks and daily rests (including split breaks, reduced rests, and extended driving allowances, together with the rule that a sufficiently long charging stop can count as a break), so as to minimize route duration while respecting the daily provisions of EU HoS regulation, battery limits, and customer service windows.

Because decisions are taken sequentially while uncertainty is revealed leg by leg, the central methodological question is about the decision architecture: when should each decision be committed, and what information should it use? We organize our solution methods along exactly this axis. Three offline methods that commit a plan before departure: a box robust plan over an interval uncertainty set, a budgeted robust plan that protects only against a bounded number of simultaneous worst-case legs, and a two-stage stochastic plan whose activity structure is fixed but whose durations adapt online, and two online architectures that re-decide at every stop: a scenario-based rolling-horizon look-ahead policy and a myopic greedy policy representing typical driver behavior. All methods are evaluated in a common discrete-event simulation, and benchmarked against a perfect-hindsight oracle, so that the value of \DEL{information and the value of} adaptivity can be measured directly.

\begin{afix}
``the value of information'' is deleted because the VSS/EVPI decomposition that
would deliver it has not been run (Section~\ref{sec:results}). ``The value of
adaptivity'' \emph{is} delivered, by the base-case comparison. Restore the full
phrase only if the VSS harness is run.
\end{afix}

Our contributions are as follows.
First, we model the full \DEL{weekly}\FIX{daily} provisions of the EU Hours-of-Service (HoS) regulation in the context of electric trucks, incorporating partial charging along a non-linear charging curve. We formulate this as a MILP and validate it against an independent exact method for the HoS rules\FIX{, reproducing the published optimum on 39 of 40 benchmark instances and improving on the 40th}.
Second, we analyze the impact of travel-time uncertainty through an extensive set of methods spanning different decision architectures, quantifying the effect of committing to a schedule upfront versus adapting en route as new information is revealed.
Third, we quantify the electrification penalty of long-haul routes by decomposing it into charging time, additional station-access time, and \DEL{wait time}\FIX{the break time that charging discharges}. This corroborates the finding of \citet{brockmann_break-and-charge_2025} that mandatory breaks provide sufficient charging windows most of the time, with charge/break coupling of up to 80\%\FIX{, while showing that the access overheads of a fast-charging stop, which that work does not price, absorb about half of the resulting credit}.
Finally, we derive policy insights regarding charging-infrastructure investment along long-haul corridors. In particular, that charger power outweighs charger density by a factor of roughly four in its effect on route duration.

\begin{acode}
\textbf{``weekly'' must become ``daily''.} \texttt{MILP.py} lines 1241--1250 put
weekly accounting explicitly out of scope: no weekly rest, no 56\,h weekly
driving cap, and the 60\,h working week is not enforced (you have already removed
that constraint from Section~\ref{sec:model} --- good). This is the first
sentence of the contribution list, so it is the first claim a referee checks, and
\citet{garaix_label-setting_2024} already own the weekly TDSP. Daily-only is
still a first in this combination.
\end{acode}

\begin{acmt}
Two smaller points. (i) ``wait time'' does not exist in the decomposition: idle
waiting is fixed to zero and Table~\ref{tab:diesel-decomp} reads $0.00$ in every
class. The three real non-charging terms are queueing, manoeuvring and bay
repositioning, and the interesting term is the negative one. (ii) ``corroborates''
undersells you --- your own decomposition shows the access overheads eat about half
the coupling credit, which is a \emph{qualification} of the closest prior work
and worth more to a referee than an agreement with it. The added clause keeps
your friendlier framing while claiming it.
\end{acmt}

\begin{atodo}
``coupling of up to 80\%'' --- the class means are \prov{68--76\,\%}
(Table~\ref{tab:diesel}) and 80\,\% is the maximum over time-window cells. Say
which. Likewise ``a factor of roughly four'': from
Section~\ref{sec:sensitivity} the ratio is \prov{4.3} against \prov{1.2--1.5},
so three to four; ``roughly four'' is the optimistic end.
\end{atodo}

The remainder of the paper is organized as follows. Section \ref{sec:litreview} reviews related work. Section \ref{sec:problem_definition} defines the problem and states the modeling assumptions. Section \ref{sec:model} presents the deterministic MILP. Section \ref{sec:method} describes the simulation framework and the {six} decision architectures. Section \ref{sec:exp} details the experimental design, Section \ref{sec:results} reports the computational results\FIX{, Section~\ref{sec:discussion} discusses their implications}. Section \ref{section:conclusion} concludes.

\begin{acmt}
The section numbering in the compiled v8 read ``Section 6 describes the
simulation framework'' because the complexity draft occupies Section 5. Once that
section is removed (Section~\ref{sec:complexity}) the numbering returns to
2--3--4--5--6--7, and this sentence is correct as written.
\end{acmt}

\section{Literature review} \label{sec:litreview}
\subsection{The truck driver scheduling problem}
Hours-of-service (HoS) regulations impose constraints on the working and driving schedule of commercial vehicle drivers. In the European Union, Regulation~(EC) No~561/2006 \citep{noauthor_regulation_2006} specifies maximum consecutive and daily driving times, mandatory break structure and minimum daily rest periods. A similar directive (2002/15/EC) limits total shift working time and restricts work during night hours. Together these rules define the truck driver scheduling problem (TDSP): given a fixed sequence of customer stops, find a feasible and cost-efficient driver schedule complying with all regulatory constraints.

The first work to model government HoS constraints in a scheduling context is \citet{xu_solving_2003}, who study a rich pickup-and-delivery problem under US rules and conjecture that finding a minimum-cost compliant schedule is NP-hard with multiple time windows. \citet{archetti_trip_2009} show that the US TDSP with single time windows is solvable in $O(n^3)$ time. For European rules, \citet{drexl_labelling_2010} formalize EU HoS as resource constraints in the elementary shortest path problem with resource constraints (ESPPRC), develop exact and heuristic labeling algorithms, and conjecture NP-completeness. \citet{goel_vehicle_2009} proposes a depth-first label-setting procedure for EU rules, which is extended in \citet{goel_truck_2010} to cover extended driving allowances and reduced rest exceptions. \citet{kok_dynamic_2010} present a dynamic programming heuristic for combined vehicle routing and driver scheduling under full EU social legislation, showing that the flexibility of complex break rules yields substantial cost reductions. \citet{prescott-gagnon_european_2010} embed EU driver rules as resources in a column-generation framework for the VRP with time windows, treating the TDSP as a pricing subproblem. \citet{rancourt_long-haul_2013} address long-haul routing under US HoS with tabu search, confirming that the choice of embedded scheduling algorithm significantly impacts solution quality. \citet{goel_hours_2014} compare EU, US, Canadian and Australian HoS regulations through optimization, finding that EU rules achieve the highest road safety at a moderate cost premium. Exact branch-and-price algorithms for the combined VRTDSP are developed by \citet{goel_exact_2017}, which is the first with proven optimality guarantees. This method is subsequently accelerated by \citet{tilk_bidirectional_2020} via bidirectional labeling. The most comprehensive EU treatment to date is \citet{garaix_label-setting_2024}, who develop a label-setting algorithm for the weekly TDSP and are the first to model the night-work rule, which limits shift working time to ten hours when night work is performed. Across this literature the driver schedule is computed once, against deterministic travel times. To our knowledge no exact EU TDSP formulation has been integrated in a sequential decision framework in which uncertain leg durations are revealed as the route is executed, which is the setting of the present paper.

\subsection{The electric vehicle routing problem}
The electric vehicle routing problem (EVRP) extends classical vehicle routing to account for the limited driving range of battery electric vehicles and the need for en-route recharging at charging stations. For a comprehensive survey we refer to \citet{erdelic_survey_2019}. The foundational formulation is the green VRP of \citet{erdogan_green_2012}, which routes alternative-fuel vehicles through refueling stations with full recharging. Exact branch-price-and-cut algorithms for the EVRPTW, covering four variants from single full recharges to multiple partial recharges, are developed by \citet{desaulniers_exact_2016}, who establish that allowing both multiple and partial recharges reduces routing cost and fleet size.

Partial recharging is studied by \citet{felipe_heuristic_2014} for a multi-technology green VRP and it is shown by \citet{keskin_partial_2016} to substantially reduce route duration . Depot-level charge scheduling over multiple days, accounting for time-dependent energy prices and battery degradation, is studied by \citet{pelletier_charge_2018}. The charging function for battery follows a concave nonlinear profile due to the constant-current/constant-voltage physics of lithium-ion cells. While most work simplify this assumption by assuming a linear charging function, it may also be linearized through the use of a piecewise-linear (PWL) approximation \citep{montoya_electric_2017}.

Uncertainty at charging stations has received growing attention: \citet{keskin_simulation-based_2021} address the EVRPTW with stochastic waiting times due to limited charger availability, showing that queuing uncertainty significantly affects routing decisions. Energy consumption uncertainty is addressed by \citet{pelletier_electric_2019} using a robust optimization framework, demonstrating that robust solutions come at only a modest cost increase. The bulk of EVRP research targets light-duty vehicles in urban or regional settings. Battery electric heavy-duty trucks (BEHDTs) differ qualitatively: much larger energy demands per leg, reliance on high-power fast charging, and binding HoS constraints that govern both when the driver may be on duty and how charging time interacts with mandatory break requirements.

\subsection{Combining TDSP and EVRP}
Despite the practical importance of long-haul BET operations, the literature combining HoS regulation with battery energy management in a single optimization model is sparse. \citet{bai_rollout-based_2023} study the problem of jointly determining optimal charging and rest decisions for an electric truck on a fixed pre-planned route under HoS constraints. They formulate a bilinear MIP minimizing total cost (charging plus rest time) and develop a rollout-based approximate solver that scales linearly with the number of stations, validated on Swedish road network data.

\citet{brockmann_break-and-charge_2025} develop a mathematical programming model that synchronizes driver break times with truck charging events under EU HoS, validated on real-world data from northwest Germany. Their key finding is that mandatory breaks create natural charging windows sufficient to match diesel route durations\FIX{; Section~\ref{sec:diesel} revisits that finding with an explicit accounting of the station-access overheads it does not price}. \citet{wan_long-haul_2024} study long-haul truck charging planning with time and energy flexibility, proposing a mixed-integer model for opportunity charging at routine stops. \citet{su_optimizing_2025} introduce a stage-expanded graph for minimizing the carbon footprint of long-haul heavy-duty e-truck trips, simultaneously optimizing path, speed, and charging, and show that the optimized charging strategy compounds the emissions benefit of electrification. \citet{vidotte_plaza_charger_2026} are the first to jointly optimize charging infrastructure siting and HoS-compliant driver scheduling for long-haul eHDT fleets, though their focus is on strategic network design rather than operational online scheduling. Finally, \citet{niyomphon_stochastic_2026} propose a chance-constrained stochastic MILP for electric freight on fixed long-haul routes with partial recharging and uncertain freight demand, but without HoS constraints.

\subsection{Uncertainty in long-haul freight transport} \label{sec:uncertaintylit}
Travel time uncertainty in long-haul road freight creates two risks: a HoS violation when accumulated driving or working time exceeds its limit before a scheduled break, and battery depletion when a leg \FIX{is driven faster than nominal and therefore} uses more energy than planned. While stochastic EVRP models have addressed energy consumption uncertainty \citep{pelletier_electric_2019} and queuing uncertainty at chargers \citep{keskin_simulation-based_2021}, travel time uncertainty and its joint effect on HoS compliance and SOC feasibility remains largely unexplored.

\begin{acode}
The neutral wording you now have is no longer wrong, but it hides the mechanism.
Under the ECR curve calibrated in \texttt{src/settings.py},
$\mathrm{ECR}(v) = A/v + B + Cv^2$ has its \emph{minimum} at $v \approx 61$\,km/h
and is increasing over the whole realizable band $[50,90]$\,km/h:

\medskip
\begin{center}\footnotesize
\begin{tabular}{lrrrr}
\toprule
$\xi$ & 0.889 (fastest) & 1.00 (nominal) & 1.20 & 1.60 (slowest)\\
speed (km/h) & 90 & 80 & 67 & 50\\
energy, 25\,km leg (kWh) & \textbf{29.40} & 27.49 & 26.04 & 26.67\\
\bottomrule
\end{tabular}
\end{center}

\medskip
So a \emph{delayed} leg consumes \emph{less} energy than planned. Stranding risk
comes from the fast tail, HoS risk from the slow tail. Section~\ref{sec:robu}
already relies on this (two independent budgets); Sections~\ref{intro} and
\ref{sec:ro} did not say it. Note also that the largest possible energy deviation
over the whole support is only $+7\,\%$ --- which is what makes 2SP's
\prov{15.5\,\%} stranding rate a statement about its \emph{plan structure} rather
than about energy noise. Worth one sentence in Section~\ref{sec:discussion}.
\end{acode}

The framework of this type of problem where decisions have to be made sequentially, and adjusted based on realized uncertainty, fits well with the rolling horizon methodology, especially when the multi stage optimization problem would be too complex to solve. The rolling-horizon paradigm is a well-established technique for multi-period optimization under gradual information revelation \citep{glomb_rolling-horizon_2022}. The broader literature on dynamic vehicle routing \citep{pillac_review_2013} covers stochastic programming, robust optimization, and online policies, and classifies our problem as one with a fully known route structure but sequentially revealed stochastic leg costs.

\subsection{Contribution}

Table~\ref{tab:litreview} positions the present work against the most closely related papers across \DEL{eight}\FIX{six} dimensions. This paper is the first to jointly model full EU HoS (including split-break sequencing, shift working-time limits, multiple rest types, and charging-break interactions), partial fast-charging with a nonlinear PWL charging curve, multiple customer stops with soft time windows on a fixed long-haul route, and stochastic travel times.

\DEL{CHECK AGAIN TABLE}

\begin{atodo}
The table has six comparison columns, not eight (corrected above). More
importantly, \textbf{it has no column for what the paper actually contributes.}
The manuscript claims the decision-architecture comparison as its methodological
contribution, but the table has no ``uncertainty handling'' and no ``benchmark''
dimension, so a referee reading only the table concludes the contribution is
``EU HoS + non-linear partial charging'' --- which
\citet{brockmann_break-and-charge_2025} largely has. Two extra columns would fix
it (``Uncert.\ handling'': deterministic / robust / stochastic / online;
``Bound'': exact / hindsight / none). Not applied here, since you asked to keep
the current structure.

\medskip
Second: the rolling-horizon and dynamic-VRP strand cited above has no row at all,
so the new column would have nothing to contrast against. Consider adding one
representative entry.
\end{atodo}

\begin{table}[htbp]
\centering
\caption{Positioning of the present work in the literature.}
\label{tab:litreview}
\scriptsize
\setlength{\tabcolsep}{3.5pt}
\resizebox{\textwidth}{!}{%
\begin{tabular}{p{3.8cm} c p{1.8cm} p{1.5cm} c c c}
\toprule
\textbf{Reference} & \textbf{HoS} & \textbf{Charging} & \textbf{Route length} & \textbf{Fixed} & \textbf{TW} & \textbf{Uncertainty}\\
\midrule
\citealt{goel_exact_2017}   & EU & $-$ & Long-haul & $-$ & Hard & $-$\\
\citealt{garaix_label-setting_2024}   & EU & $-$ & Long-haul & $\checkmark$ & Hard & $-$\\
\citealt{montoya_electric_2017}   & $-$ & Partial + NL & Regional & $-$ & Hard & $-$\\
\citealt{keskin_simulation-based_2021}   & $-$ & Partial + L & Long-haul & $-$ & Hard & CS queues\\
\citealt{niyomphon_stochastic_2026}   & $-$ & Partial + L & Long-haul & $-$ & $-$ & Demand \\
\citealt{bai_rollout-based_2023}   & EU & Partial$+$L & Day & $\checkmark$ & $-$ & $-$\\
\citealt{brockmann_break-and-charge_2025}   & EU & Partial$+$NL  & Day & $\checkmark$ & $-$ & $-$ \\
\citealt{vidotte_plaza_charger_2026}   & Brazil & Full$+$L &  Long-haul & $-$ & Hard & $-$ \\
\midrule
This paper & {EU} & {Partial$+$NL} & {Long-haul} & $\checkmark$ & {Soft} & Travel time\\
\bottomrule
\end{tabular}}
\par\smallskip
\begin{minipage}{\textwidth}
\footnotesize
\textbf{Charging}: L/NL: linear/non-linear charging function.
\textbf{TW}: customer time windows.
\end{minipage}
\end{table}
